\subsection{Cross-Dataset Generalisation}
\label{sec:ood}

To assess whether representations learnt on DHCD transfer beyond the training
distribution, we evaluate both \modelname{} and the reproduced MallaNet
baseline on two independently collected Devanagari datasets covering different
writers, scanner workflows, and collection years.
No fine-tuning or parameter adaptation is applied in either case.

\paragraph{NHCD (Pant 2012).}
The Nepali Handwritten Character Dataset~\cite{pant2012nhcd} was assembled at
Tribhuvan University, Nepal, in 2012 as part of an independent research
programme---three years before DHCD and with no overlap in
writers (40 named contributors, identifiers encoded in filenames).
It covers the same 46-class inventory (36 consonants\,+\,10 digits) in
$28{\times}28$ JPEG images with dark-ink-on-white polarity.
Before evaluation we polarity-invert and bilinearly resize to $32{\times}32$
to match the DHCD convention; no labels are used.

Table~\ref{tab:ood} summarises the outcome.
\modelname{} achieves \textbf{78.92\%} zero-shot overall
(95.80\% on digits, 72.33\% on consonants).
MallaNet collapses to \textbf{23.99\%}
(58.85\% on digits, 10.38\% on consonants)---a \textbf{55\,pp gap}
despite identical in-distribution performance on DHCD.
The disparity on consonants (+62\,pp) is the most striking:
MallaNet's homogeneous filter capsule (HFC) layers and fixed
$(0.5,\,0.5)$ normalisation appear to overfit to DHCD's specific
stroke rendering, while the distilled student's soft-label training
from a diverse teacher ensemble produces representations that are
less dataset-specific.

One systematic error cluster is linguistically informative.
For both models, the four dental/retroflex consonant pairs
(ta/\d{t}a, tha/\d{t}ha, da/\d{d}a, dha/\d{d}ha) produce
near-zero accuracy with systematic cross-class confusion.
These pairs are visually distinguished by a single diacritic stroke
whose rendering conventions differ between Indian and Nepali handwriting
traditions---a script-level domain gap that neither model resolves,
but that \modelname{} handles better everywhere else.

\paragraph{Prashanth et al.\ 2021.}
Prashanth et al.~\cite{prashanth2021} provide 22{,}500 handwritten
Devanagari digit images (10 classes, 2{,}250 per digit, $28{\times}28$
JPEG) collected independently from DHCD, with matching dark-background
polarity.
After $32{\times}32$ resize, \modelname{} achieves \textbf{79.31\%}
zero-shot; MallaNet achieves \textbf{72.60\%}.
Both results are consistent with the NHCD digit outcome and with the
CMATERdb transfer result reported in Section~\ref{sec:transfer},
confirming that the digit representation gap is reproducible across
independent datasets.
Devanagari digit~5 confuses with digit~7 at high rates for both
models---a genuine visual similarity in the script rather than a
model-specific artefact.

\begin{table}[t]\centering
\caption{Zero-shot cross-dataset accuracy. Both models trained on DHCD
  only; no fine-tuning on the target dataset. NHCD polarity-inverted
  before evaluation; Prashanth polarity matches DHCD.
  $n_{\text{NHCD}}=10{,}260$;
  $n_{\text{Prashanth}}=22{,}500$ (digit classes only).}
\label{tab:ood}
\begin{tabular}{lcccc}\toprule
& \multicolumn{3}{c}{NHCD (Pant 2012, 46 classes)} & Prashanth 2021 \\
\cmidrule(lr){2-4}
Model & Overall & Consonants & Digits & (10 digit classes) \\\midrule
\modelname{} (ours)      & \textbf{78.92\%} & \textbf{72.33\%} & \textbf{95.80\%} & \textbf{79.31\%} \\
MallaNet~\cite{malla2025} & 23.99\%         & 10.38\%          & 58.85\%          & 72.60\%          \\
\midrule
Gap                       & $+$54.9\,pp     & $+$61.9\,pp      & $+$36.9\,pp      & $+$6.7\,pp       \\\bottomrule
\end{tabular}\end{table}

These results reveal that matched in-distribution accuracy conceals a
substantial generalisation gap.
The compact distilled model transfers reliably across writer pools and
scanner workflows; the larger model does not.
This finding echoes the corruption-robustness result of
Section~\ref{sec:robustness}: both out-of-distribution evaluations
consistently favour \modelname{}, suggesting that
knowledge distillation from a diverse teacher ensemble provides a
generalisation benefit that raw accuracy on a saturated benchmark
cannot detect.
